% from ega2-1-fr.tex
\oldpage[II]{183}
la même hypothèse lorsque l'on suppose $X$ noethérien, car alors $V$ est aussi
noethérien et l'assertion découle cette fois de
(\hyperref[I.6.3.5-fr]{I, 6.3.5}).

\begin{corollary}[8.9.2]
\label{II.8.9.2-fr}
Si le morphisme $p:X\to Y$ est propre, on peut dans l'énoncé de
(\hyperref[II.8.9.1-fr]{8.9.1}), supposer $q$ propre, et remplacer
«~immersion ouverte~» par «~isomorphisme~».
\end{corollary}

De façon imagée, on peut dire (lorsque $p:X\to Y$ est propre) que
$\mathcal{L}$ est ample relativement à $p$ si et seulement si on peut
«~contracter~» la section nulle du fibré vectoriel $\mathbf{V}(\mathcal{L})$
en le préschéma de base $Y$. Un cas particulier important est celui où $Y$
est le spectre d'un corps, et où l'opération de «~contraction~» consiste donc
à contracter la section nulle de $\mathbf{V}(\mathcal{L})$ \emph{en un seul
point}.

\begin{env}[8.9.3]
\label{II.8.9.3-fr}
La nécessité des conditions de (\hyperref[II.8.9.1-fr]{8.9.1}) et le
corollaire (\hyperref[II.8.9.2-fr]{8.9.2}) découlent aussitôt de
(\hyperref[II.8.8.2-fr]{8.8.2}) et
(\hyperref[II.8.8.4-fr]{8.8.4}).

Pour démontrer la suffisance de (\hyperref[II.8.9.1-fr]{8.9.1}), nous
considérerons une situation un peu plus générale. Pour cela, posons (avec les
notations de (\hyperref[II.8.8.2-fr]{8.8.2}))
\[
  \mathcal{S}'=\bigoplus_{n\geqslant 0}\mathcal{L}^{\otimes n}
\]
et
\[
  V=\mathbf{V}(\mathcal{L})=\operatorname{Spec}(\mathcal{S}').
\]
Le sous-préschéma fermé $j(X)$, section nulle de $\mathbf{V}(\mathcal{L})$,
est défini par le faisceau quasi-cohérent d'idéaux
$\mathcal{J}=(\mathcal{S}'_+)^{\sim}$ de $\mathcal{O}_V$
(\hyperref[II.1.4.10-fr]{1.4.10}). Ce $\mathcal{O}_V$-Module est
\emph{inversible}, car la question est locale sur $X$ et cela revient à
remarquer que l'idéal $TA[T]$ dans un anneau de polynômes $A[T]$ est un
$A[T]$-module libre monogène. En outre, il est immédiat (toujours parce que la
question est locale sur $X$) que
\[
  \mathcal{L}=j^*(\mathcal{J})
\]
et
\[
  j_*(\mathcal{L})=\mathcal{J}/\mathcal{J}^2.
\]
D'autre part, si
\[
  \pi:\mathbf{V}(\mathcal{L})\longrightarrow X
\]
est le morphisme structural, on a
$\pi_*(\mathcal{J})=\mathcal{S}'_+$ et
$\pi_*(\mathcal{J}/\mathcal{J}^2)=\mathcal{L}$~; on a donc deux
homomorphismes canoniques
$\mathcal{L}\to\pi_*(\mathcal{J})\to\mathcal{L}$, le premier étant
l'injection canonique $\mathcal{L}\to\mathcal{S}'_+$, le second la projection
canonique de $\mathcal{S}'_+$ sur $\mathcal{S}'_1=\mathcal{L}$, et leur
composé est l'identité. On peut d'ailleurs plonger canoniquement
\[
  \pi_*(\mathcal{J})=\mathcal{S}'_+
  =\bigoplus_{n\geqslant 1}\mathcal{L}^{\otimes n}
\]
dans le \emph{produit}
\[
  \prod_{n\geqslant 1}\mathcal{L}^{\otimes n}
  =\varprojlim_n\pi_*(\mathcal{J}/\mathcal{J}^{n+1})
\]
(puisque
$\pi_*(\mathcal{J}/\mathcal{J}^{n+1})=
\mathcal{L}\oplus\mathcal{L}^{\otimes2}\oplus\cdots\oplus
\mathcal{L}^{\otimes n}$), et on a donc encore deux homomorphismes canoniques
\begin{equation}
\label{II.8.9.3.1-fr}
  \mathcal{L}\longrightarrow
  \varprojlim_n\pi_*(\mathcal{J}/\mathcal{J}^{n+1})
  \longrightarrow\mathcal{L}
\tag{8.9.3.1}
\end{equation}
dont le composé est l'identité.

Cela étant, la généralisation de (\hyperref[II.8.9.1-fr]{8.9.1}) que nous
allons prouver est la suivante~:
\end{env}

\begin{proposition}[8.9.4]
\label{II.8.9.4-fr}
Soient $Y$ un préschéma, $V$ un $Y$-préschéma, $X$ un sous-préschéma fermé de
$V$, défini par un Idéal $\mathcal{J}$ de $\mathcal{O}_V$, qui est un
$\mathcal{O}_V$-Module inversible~; si $j:X\to V$ est

